
\section{Worked example: hyperfine splitting}
\textbf{Question type.} Given \(I=\frac52\), \(J=\frac12\) and magnetic-dipole
hyperfine constant \(A\), determine the allowed \(F\), energies,
degeneracies and separation.

\textbf{Step 1: allowed total angular momenta.}
\[
F=|I-J|,|I-J|+1,\ldots,I+J
\quad\Longrightarrow\quad F=2,\ 3.
\]
\textbf{Step 2: derive rather than blindly apply the formula.}
\[
\hat{\mathbf F}^2=\hat{\mathbf I}^2+\hat{\mathbf J}^2+
2\hat{\mathbf I}\cdot\hat{\mathbf J},
\quad
\hat H=\frac A{\hbar^2}\hat{\mathbf I}\cdot\hat{\mathbf J},
\]
so
\[
E_F=\frac A2[F(F+1)-I(I+1)-J(J+1)].
\]
\textbf{Step 3: substitute.} Here \(I(I+1)=35/4\),
\(J(J+1)=3/4\).
\begin{align*}
E_{F=2}&=\frac A2\left(6-\frac{35}{4}-\frac34\right)
=-\frac{7A}{4},\\
E_{F=3}&=\frac A2\left(12-\frac{35}{4}-\frac34\right)
=\frac{5A}{4},\\
E_3-E_2&=3A.
\end{align*}
\textbf{Step 4: degeneracies and check.}
\(2F+1=5\) or \(7\) magnetic substates.
Their sum \(12=(2I+1)(2J+1)\) verifies completeness.

\begin{methodbox}
\textbf{Reusable exam method.} Couple \(I,J\) to list \(F\);
use \(F(F+1)\); compute each energy; count \(2F+1\);
check the dimension and sign of \(A\).
\end{methodbox}
\begin{cautionbox}
Do not replace the eigenvalue \(F(F+1)\hbar^2\) with \(F^2\).
If \(A<0\), the ordering reverses. The degeneracy \(2F+1\)
is \emph{not} the number of distinct hyperfine energies.
\end{cautionbox}
